\labeledsection{Hidden Markov Models}{sec:hmm}
\titledbox{Motivation}{
\begin{itemize}
\item The world changes in \textit{stochastically predictable ways}
\item For example, the weather changes, but the weather tomorrow is somewhat predictable given today's weather and other factors.
\item The stock market changes, but the stock price tomorrow is probably related to today's price 
\item A patient's blood sugar changes, but their blood sugar is related to their blood sugar 10 minutes ago 
\end{itemize}
}

\questionbox{
How can we model this?
}

\anondefinition{
Given a \linkterm{probability space}{def:prob_space} $\tuple{\Omega, \mathcal{F}, P}$, and a time structure $S$ with an order $\tuple{S, \preceq}$. Then a \linkterm{sequence}{sequence} of \linkterm{random variables}{def:random_variable} $(X_t)t \in S$ with $\text{dom}(X_t) = D$ is called \defemph{stochastic process} over the \defemph{time strcuture} $S$
}{stochastic_process}

\commandnote{
$X_t$ models the outcome of the \linkterm{random variable}{def:random_variable} $X$ at a given time step $t$. The sample sapce $\Omega$ corresponds to the \linkterm{set}{def:set} of all possible \linkterm{sequences}{sequence} of outcomes.
}

\redbox{
We will almost exclusively use $\tuple{S, \preceq} = \tuple{\mathbb{N}, \leq}$
}

\notationbox{
Given a \linkterm{stochastic process}{stochastic_process} $X_t$ over $S$ and $a,b \in S$ with $a \preceq b$, we write $X_{a:b}$ for the \linkterm{sequence}{sequence} $X_a, X_{a+1}, \cdots, X_{b-1}, X_b$ and $X_{a:b}^{=x}$ for $X_a = x_a, \cdots, X_b = x_b$
}





\ideabox{
Construct a \linkterm{Bayesian Network}{def:bayes_net} from these variables
}

\anondefinition{
Let $(X_t)_{t \in S}$ be a \linkterm{stochastic process}{stochastic_process}. We say that $X$ has \defemph{n-th order Markov property} iff $X_t$ only depends on a bounded subset of $X_{0:t-1}$, i.e. for all $t \in S$ we have $\probd{X_t \mid X_0, \cdots, X_{t-1}} = \probd{X_t \mid X_{t-n}, \cdots , X_{t-1}}$ for some $n \in S$.

A \linkterm{stochastic process}{stochastic_process} with the \textbf{Markov} property for some $n$ is called an \defemph{n-th order Markov process}
}{markov_property}

\definition{First-Order Markov Property (Markov Chain)}{
\[
\probd{X_t \mid X_{0:t-1}} = \probd{X_t \mid X_{t-1}}
\]

\begin{figure}[H]
\centering
\begin{tikzpicture}[node distance=1.5cm]
  % Left ellipsis
  \node at (-0.75, 0) {$\cdots$};
  
  % Nodes
  \node[vertex] (xt_2) at (0, 0) {$X_{t-2}$};
  \node[vertex] (xt_1) at (1.5, 0) {$X_{t-1}$};
  \node[vertex] (xt) at (3, 0) {$X_t$};
  \node[vertex] (xt1) at (4.5, 0) {$X_{t+1}$};
  \node[vertex] (xt2) at (6, 0) {$X_{t+2}$};
  
  % Arrows between nodes
  \draw[arc] (xt_2) -- (xt_1);
  \draw[arc] (xt_1) -- (xt);
  \draw[arc] (xt) -- (xt1);
  \draw[arc] (xt1) -- (xt2);
  
  % Right ellipsis
  \node at (6.75, 0) {$\cdots$};
\end{tikzpicture}
\end{figure}

}{markov_chain}

\definition{Second-Order Markov Property}{
\[
\probd{X_t \mid X_{0:t-1}} = \probd{X_t \mid X_{t-2}, X_{t-1}}
\]

\begin{figure}[H]
\centering
\begin{tikzpicture}[node distance=1.5cm]
  % Left ellipsis
  \node at (-0.75, 0) {$\cdots$};
  
  % Nodes
  \node[vertex] (xt_2) at (0, 0) {$X_{t-2}$};
  \node[vertex] (xt_1) at (1.5, 0) {$X_{t-1}$};
  \node[vertex] (xt) at (3, 0) {$X_t$};
  \node[vertex] (xt1) at (4.5, 0) {$X_{t+1}$};
  \node[vertex] (xt2) at (6, 0) {$X_{t+2}$};
  
  % Main arrows between nodes
  \draw[arc] (xt_2) -- (xt_1);
  \draw[arc] (xt_1) -- (xt);
  \draw[arc] (xt) -- (xt1);
  \draw[arc] (xt1) -- (xt2);
  
  % Arced arrows showing second-order dependencies (above)
  \draw[arc, bend left=60] (xt_2) to (xt);
  \draw[arc, bend left=60] (xt_1) to (xt1);
  \draw[arc, bend left=60] (xt) to (xt2);
  
  % Right ellipsis
  \node at (6.75, 0) {$\cdots$};
\end{tikzpicture}
\end{figure}
}{}

\commandnote{
Given a \linkterm{Markov chain}{markov_chain}, and we are now at time $t$, then the future is \linkterm{conditionally independent}{def:cond_independence} of the past given the present: 
\[
\probd{X_{t+1} \mid X_{0:t}} = \probd{X_{t+1} \mid X_t}
\]
}

\anondefinition{
Let $(X_t)_{t \in S}$ be a \linkterm{stochastic process}{stochastic_process} where $X$ has the \linkterm{Markov property}{markov_property} (for any $n$). We call the \linkterm{probability distribution}{def:prob_distribution} $\probd{X_t \mid X_{0:t-1}}$ the \defemph{transition model} of the \linkterm{stochastic process}{stochastic_process} $(X_t)_{t \in S}$
}{transition_model_hmm}

\definition{Stationary Assumption}{
A \linkterm{Markov chain}{markov_chain} is called \defemph{stationary} if the \linkterm{transition model}{transition_model_hmm} is \linkterm{independent}{def:independence} of time, i.e. $\probd{X_t \mid X_{t-1}}$ is the same for all $t$
}{stationary_hmm}

\Example{
Consider we represent the weather as a \linkterm{Markov chain}{markov_chain}, we introduce a variable $X$ with domain $\set{\text{sun}, \text{rain}}$ and we fix the initial distribution $\probd{X} = \tuple{1, 0}$ (sunny), and we have a \linkterm{stationary}{stationary_hmm} \linkterm{transition model}{transition_model_hmm} that might be: 

\begin{tabular}{|c|c|c|} \hline
    $X_{t-1}$ & $X_t$ & $P(X_t \mid X_{t-1})$ \\ \hline
    sun & sun & 0.9 \\ \hline
    sun & rain & 0.1 \\ \hline
    rain & sun & 0.3 \\ \hline
    rain & rain & 0.7 \\ \hline
\end{tabular}
}{weather_example_hmm}

\titledbox{Different ways of representing transition models}{
If a \linkterm{Markov chain}{markov_chain} $X$ is \linkterm{stationary}{stationary_hmm} and with \linkterm{finite}{set_cardinality} \linkterm{domain}{domain_of_discourse}, we can represent the \linkterm{transition model}{transition_model_hmm} as a matrix $T$ where $T_{ij} := P(X_{t=j} \mid X_{t-1=i})$ 

\[
T =
\begin{pmatrix}
P(\text{sun}\mid \text{sun}) & P(\text{rain}\mid \text{sun}) \\
P(\text{sun}\mid \text{rain}) & P(\text{rain}\mid \text{rain})
\end{pmatrix}
\] 

\[
\begin{pmatrix}
0.9 & 0.1 \\
0.3 & 0.7
\end{pmatrix}.
\]

We can also represent it to look like a finite state machine :
\begin{figure}[H]
\centering
\begin{tikzpicture}[node distance=3.5cm, vertex/.style={ellipse, draw, minimum width=52pt, minimum height=30pt, inner sep=0pt}]
  \node[vertex] (sun) at (0, 0) {sun};
  \node[vertex] (rain) at (4, 0) {rain};

  % Initial state
  \draw[arc] (-2, 0) -- (sun);

  % Transitions
  \draw[arc] (sun) edge[loop above] node {$0.9$} (sun);
  \draw[arc] (rain) edge[loop above] node {$0.7$} (rain);
  \draw[arc] (sun) to[bend left=18] node[above] {$0.1$} (rain);
  \draw[arc] (rain) to[bend left=18] node[below] {$0.3$} (sun);
\end{tikzpicture}
\end{figure}
}

\questionbox{
Given our weather example (\exampleref{weather_example_hmm}), what is the probability distribution of $X$ after one time step (knowing it is initially sunny)?
}

\answerbox{
\begin{itemize}
    \item Trivially from the transition model $P(X_2 = sun \mid X_1 = sun) = 0.9$. Hence, $\probd{X_2} = \tuple{0.9, 0.1}$ But let's try to reason about it.
    \item We know that for $t=1$, $P(\text{sun}) = 1$ and $P(\text{rain}) = 0$, i.e. the initial distribution is $\probd{X_1} = \tuple{1, 0}$. We also have a \linkterm{stationary}{stationary_hmm} \linkterm{transition model}{transition_model_hmm}. We are interested in knowing $\probd{X_2}$.
\end{itemize}

We compute it for \textit{sun} by \linkterm{marginalization}{def:marginalization} over the different values of $X_{t=1}$:
\begin{align*}
P(X_2 = \text{sun}) &= \sum_{x \in \set{\text{sun}, \text{rain}}} P(X_2 = \text{sun}, X_1 = x) \\ 
&=\sum_{x \in \set{\text{sun}, \text{rain}}} P(X_2 = \text{sun} \mid X_1 = x) \cdot P(X_1 = x) \\
P(X_2 = \text{sun}) &= P(X_2 = \text{sun} \mid X_1 = \text{sun}) \cdot P(X_1 = \text{sun}) + 
P(X_2 = \text{sun} \mid X_1 = \text{rain}) \cdot P(X_1 = \text{rain}) \\ 
&= 0.9 \cdot 1.0 + 0.1 \cdot 0 = 0.9   
\end{align*}
}

\redbox{
What we did is the simplest algorithm for \linkterm{Markov chains}{markov_chain}, it is called \textbf{mini forward}
}

\definition{Mini Forward Algorithm}{
Given a \linkterm{stationary}{stationary_hmm} \linkterm{Markov chain}{markov_chain} $X$ with initial distribution $\probd{X_1}$, we the distribution at any time step $t$ by recursively applying:

\[
\probd{X_n} = \sum_{x_{n-1}} P(X_n \mid X_{n-1} = x_{n-1}) \cdot P(X_{n-1} = x_{n-1})
\]

for $n = 2, 3, \ldots, t$. At each step, we \linkterm{marginalize}{def:marginalization} out the previous time step by weighting each possible transition by its prior probability.
}{mini_forward}

\centeredtitledbox{Conversion}{
If we were to do this for infinitely many time steps, the \linkterm{probability distribution}{def:prob_distribution} $\probd{X}$ of the \linkterm{Markov chain}{markov_chain} $X$ converges.

\[
\probd{X_1} = \tuple{1.0 , 0.0} \leadsto \probd{X_2} = \tuple{0.9 , 0.1}
\leadsto \probd{X_3} = \tuple{0.84 , 0.16}
\leadsto 
\]

\[
\probd{X_4} = \tuple{0.804 , 0.196} \leadsto \cdots \leadsto \probd{X_{\infty}} = \tuple{0.75, 0.25}
\]

This applies to any initial distribution $\probd{X} = \tuple{p, 1-p}$!
}

\anondefinition{
As $t \to \infty$, $\probd{X}$ converges towards a fixed point called the \defemph{stationary distribution} of the \linkterm{Markov chain}{markov_chain} $X$
}{stationary_distribution_hmm}

\redbox{
In other words, the \linkterm{Markov chain}{markov_chain} essentially ``forgets'' its initial state; the long-term probabilities become completely independent of the starting conditions.
}

\titledbox{Computing the Stationary Distribution}{
To compute the \linkterm{stationary distribution}{stationary_distribution_hmm} $S$, we can use the matrix representation of the \linkterm{transition model}{transition_model_hmm} $T$. Because the stationary distribution is a fixed point, passing it through the transition model returns the same distribution.

If we treat the distribution $S$ as a column vector and $T$ is the transition matrix (as defined earlier, where rows sum to 1), then taking the next step involves the transpose of $T$. The stationary distribution satisfies:
\[
S = T^T \cdot S
\]
\textit{Note: If you were to represent $S$ as a row vector instead, the equation would simply be $S = S \cdot T$.}

To compute $S$ (as a column vector), we solve the system of linear equations:
\[
(T^T - I) \cdot S = 0
\]
along with the necessary constraint that all probabilities must sum to 1 ($\sum_i S_i = 1$). This means $S$ is an eigenvector of $T^T$ corresponding to the eigenvalue $\lambda = 1$.
}

\commandnote{
\linkterm{Markov chains}{markov_chain} are not so useful for most agents, usually we need another ingredient, observations. 

We begin to speak of \textbf{hidden Markov models} (HMMs) when we have an underlying \linkterm{Markov chain}{markov_chain} $X$ that is \textbf{hidden} to us and we \textbf{observe} some \linkterm{evidence}{def:bayesian_terms} $E$ that is affected by that \linkterm{Markov chain}{markov_chain}.
}


\anondefinition{
We call a \linkterm{stochastic process}{stochastic_process} of \textit{hidden} variables a \defemph{state variable}
}{state_variable_hmm}

\Example{
You are a security guard in a secret underground facility, want to know
it if is raining outside. Your only source of information is whether the director comes in with an umbrella.
\begin{itemize}
    \item we have a \linkterm{stochastic process}{stochastic_process} $\text{Rain}_0, \text{Rain}_1, \text{Rain}_2, \cdots$ of \defemph{hidden} variables, and 
    \item a related \linkterm{stochastic process}{stochastic_process} $\text{Umbrella}_0, \text{Umbrella}_1, \text{Umbrella}_2, \cdots$ of \linkterm{evidence}{def:bayesian_terms} variables.
\end{itemize}
And a combined \linkterm{stochastic process}{stochastic_process} $\tuple{\text{Rain}_0, \text{Umbrella}_0}, \tuple{\text{Rain}_1, \text{Umbrella}_1}, \cdots$

\textbf{Note}: $\text{Umbrella}_t$ only depends on $\text{Rain}_t$, not on e.g. $\text{Umbrella}_{t-1}$

We model this as a \linkterm{Bayesian Network}{def:bayes_net}:

\begin{figure}[H]
\centering
\begin{tikzpicture}[node distance=2.5cm, vertex/.style={ellipse, draw, minimum width=60pt, minimum height=40pt, inner sep=0pt}]
  % Left ellipsis
  \node at (-1.5, 2) {$\cdots$};
  
  % Rain nodes (top row)
  \node[vertex] (rain_t_1) at (0, 2) {$\text{Rain}_{t-1}$};
  \node[vertex] (rain_t) at (2.5, 2) {$\text{Rain}_t$};
  \node[vertex] (rain_t1) at (5, 2) {$\text{Rain}_{t+1}$};
  
  % Arrows between rain nodes
  \draw[arc] (rain_t_1) -- (rain_t);
  \draw[arc] (rain_t) -- (rain_t1);
  
  % Right ellipsis
  \node at (6.5, 2) {$\cdots$};
  
  % Umbrella nodes (bottom row)
  \node[vertex] (umb_t_1) at (0, 0) {$\text{Umbrella}_{t-1}$};
  \node[vertex] (umb_t) at (2.5, 0) {$\text{Umbrella}_t$};
  \node[vertex] (umb_t1) at (5, 0) {$\text{Umbrella}_{t+1}$};
  
  % Vertical arrows from rain to umbrella
  \draw[arc] (rain_t_1) -- (umb_t_1);
  \draw[arc] (rain_t) -- (umb_t);
  \draw[arc] (rain_t1) -- (umb_t1);
\end{tikzpicture}
\end{figure}
}{umbrella_example}

\anondefinition{
We call the \linkterm{conditional probability}{def:conditional_prob} distribution $\probd{E_t \mid X_{0:t}, E_{1:t-1}}$ a \defemph{sensor model}
}{sensor_model_hmm}

\anondefinition{
We say that a \linkterm{sensor model}{sensor_model_hmm} has the \defemph{sensor Markov property}, iff $\probd{E_t \mid X_{0:t}, E_{1:t-1}} = \probd{E_t \mid X_t}$, i.e. the \linkterm{sensor model}{sensor_model_hmm} depends only the current state.
}{sensor_markov_property}

\titledbox{Observation Matrix}{
If a \linkterm{sensor model}{sensor_model_hmm} has the \linkterm{sensor Markov property}{sensor_markov_property}, we can represent each observation $E_t  =e_t$ at time $t$ as the \textit{diagonal} matrix $O_t$ with $O_{t_{ii}} := P(E_t = e_t \mid X_t = i)$
}

\anondefinition{
A pair $\tuple{X,E}$ where $X$ is a \linkterm{stationary}{stationary_hmm} \linkterm{Markov chain}{markov_chain} (hidden state variable), and $E$ has the \linkterm{sensor markov property}{sensor_markov_property} (\linkterm{evidence}{def:bayesian_terms}) is called a (\linkterm{stationary}{stationary_hmm}) \defemph{Hidden Markov Model (HMM)}
}{hmm}

\Example{
Our umbrella example (\exampleref{umbrella_example}) is a \linkterm{HMM}{hmm} with the following example \linkterm{transition model}{transition_model_hmm} and \linkterm{sensor model}{sensor_model_hmm}:

\begin{figure}[H]
\centering
\begin{tikzpicture}[node distance=7cm, vertex/.style={ellipse, draw, minimum width=60pt, minimum height=40pt, inner sep=0pt}]
  % Left ellipsis
  \node at (-2, 4) {$\cdots$};
  
  % Rain nodes (top row)
  \node[vertex] (rain_t_1) at (0, 4) {$\text{Rain}_{t-1}$};
  \node[vertex] (rain_t) at (5, 4) {$\text{Rain}_t$};
  \node[vertex] (rain_t1) at (10, 4) {$\text{Rain}_{t+1}$};
  
  % Arrows between rain nodes
  \draw[arc] (rain_t_1) -- node[above=0.7cm] {
    \begin{tabular}{|c|c|} \hline
        $R_{t-1}$ & $P(R_t = T \mid R_{t-1})$ \\ \hline 
        $T$ & $0.7$ \\ \hline
        $F$ & $0.3$ \\ \hline
    \end{tabular}
  } (rain_t);
  \draw[arc] (rain_t) -- (rain_t1);
  
  % Right ellipsis
  \node at (12, 4) {$\cdots$};
  
  % Umbrella nodes (bottom row)
  \node[vertex] (umb_t_1) at (0, 0) {$\text{Umbrella}_{t-1}$};
  \node[vertex] (umb_t) at (5, 0) {$\text{Umbrella}_t$};
  \node[vertex] (umb_t1) at (10, 0) {$\text{Umbrella}_{t+1}$};
  
  % Vertical arrows from rain to umbrella
  \draw[arc] (rain_t_1) -- (umb_t_1);
  \draw[arc] (rain_t) -- node[right=0.2cm] {
    \begin{tabular}{|c|c|} \hline
        $R_{t}$ & $P(U_t = T \mid R_t)$ \\ \hline 
        $T$ & $0.9$ \\ \hline
        $F$ & $0.2$ \\ \hline
    \end{tabular}
  } (umb_t);
  \draw[arc] (rain_t1) -- (umb_t1);
\end{tikzpicture}
\end{figure}
}{}


\titledbox{Umbrella Matrices}{
\[
O_T = \begin{pmatrix}
0.9 & 0 \\
0 & 0.2
\end{pmatrix}
\qquad
O_F = \begin{pmatrix}
0.1 & 0 \\
0 & 0.8
\end{pmatrix}
\qquad
T = \begin{pmatrix}
0.7 & 0.3 \\
0.3 & 0.7
\end{pmatrix}
\]
}


\definition{Markov Inference Tasks}{
Gievn a \linkterm{Markov chain}{markov_chain} with state variables $X_t$ and \linkterm{evidence}{def:bayesian_terms} variables $E_t$, we are interested in the following \defemph{Markov inference tasks}:
\begin{itemize}
\item \defemph{Filtering (Monitoring)}: The task of tracking / updating the \linkterm{probability distribution}{def:prob_distribution} $B_t (X) = \probd{X_t \mid E_{1:t}^{=e}}$ (\defemph{belief state}) over time. i.e. given sequence $e_1, \cdots, e_t$ of observations up until time $t$, we compute the likely state of the world at current time $t$.
\item \defemph{Prediction (State Estimation)}: Given sequence $e_1, \cdots, e_t$ of observations up until time $t$, compute the likely \textit{future} state of the world at time $t+k$ for $k > 0$. i.e. calculating the \linkterm{probability distribution}{def:prob_distribution} $\probd{X_{t+k} \mid E_{1:t}^{=e}}$.
\item \defemph{Smoothing (Hindsight)}: Given sequence $e_1, \cdots, e_t$ of observations up until time $t$, compute the likely \textit{past} state of the world at time $t-k$ for $0 < k < t$. i.e. calculate the \linkterm{probability distribution}{def:prob_distribution} $\probd{X_{t-k} \mid E_{1:t}^{=e}}$
\item \defemph{Most Likely Explanation (MLE)}: Given sequence $e_1, \cdots, e_t$ of observations up until time $t$, copmute the most likely sequence of states that led to these observations, i.e. calculate:
\[
\argmax{x} P(X_{1:t}^{=x} \mid E_{1:t}^{=e})
\]
\end{itemize}
}{markov_inference}

\commandnote{
The most likely sequence of states is \textit{not} (necessarily) the sequence of most likely states
}


\anondefinition{
Given a \linkterm{Markov chain}{markov_chain} $X$ and \linkterm{evidence}{def:bayesian_terms} variables $E$, let $\probd{X_0}$ be the initial \linkterm{prior}{def:bayesian_terms} (initial belief), then we can compute the full \linkterm{JPD}{def:full_joint_prob_distribution} as:
\[
\probd{X_{0:t}, E_{1:t}} = \probd{X_0} \cdot \left( \prod_{i=1}^{t} \probd{X_i \mid X_{i-1}} \cdot \probd{E_i \mid X_i} \right)
\]
}{jpd_hmm}

\titledbox{Explanation}{
This is a \linkterm{Bayesian Network}{def:bayes_net} (recall \thmref{jpd_BN}):

\begin{figure}[H]
\centering
\begin{tikzpicture}[node distance=2cm, vertex/.style={ellipse, draw, minimum width=30pt, minimum height=30pt, inner sep=2pt}]
  % Nodes for X (top)
  \node[vertex, label=above:{\small initial belief}] (x0) at (0, 2) {$X_0$};
  \node[vertex] (x1) at (2, 2) {$X_1$};
  \node[vertex] (x2) at (4, 2) {$X_2$};
  \node[vertex] (x3) at (6, 2) {$X_3$};
  \node (xdots) at (8, 2) {$\cdots$};
  
  % Nodes for E (bottom)
  \node[vertex] (e1) at (2, 0) {$E_1$};
  \node[vertex] (e2) at (4, 0) {$E_2$};
  \node[vertex] (e3) at (6, 0) {$E_3$};
  \node (edots) at (8, 0) {$\cdots$};
  
  % Horizontal arrows for states
  \draw[arc] (x0) -- (x1);
  \draw[arc] (x1) -- (x2);
  \draw[arc] (x2) -- (x3);
  \draw[arc] (x3) -- (xdots);
  
  % Vertical arrows for evidence
  \draw[arc] (x1) -- (e1);
  \draw[arc] (x2) -- (e2);
  \draw[arc] (x3) -- (e3);
\end{tikzpicture}
\end{figure}
}

\titledbox{Filtering Base Case}{
We want to compute the belief state at $t=1$, which is $B(X_1) = \probd{X_1 \mid E_1}$. We can derive this in two steps as follows:

By Bayes' theorem, we incorporate the new evidence to update our belief:
\begin{align*}
\probd{X_1 \mid E_1} &= \frac{\probd{E_1 \mid X_1} \cdot \probd{X_1}}{\probd{E_1}} \\
&= \alpha (\probd{E_1 \mid X_1} \cdot \probd{X_1})
\end{align*}
Here $\probd{E_1 \mid X_1}$ is our \linkterm{sensor model}{sensor_model_hmm}.

Next, we expand the prior prediction $\probd{X_1}$ by \linkterm{marginalizing}{def:marginalization} over the initial state $X_0$:
\begin{align*}
\probd{X_1} &= \sum_{x_0} \probd{X_1, X_0 = x_0} \\
&= \sum_{x_0} \probd{X_1 \mid X_0 = x_0} \cdot \probd{X_0 = x_0}
\end{align*}
Here, $\probd{X_1 \mid X_0}$ is the \linkterm{transition model}{transition_model_hmm}, and $\probd{X_0}$ is the initial belief prior.

Combining these gives the complete filtering formula for $t=1$:
\[
\probd{X_1 \mid E_1} = \alpha \left( \probd{E_1 \mid X_1} \sum_{x_0} \probd{X_1 \mid x_0} \cdot \probd{x_0} \right)
\]
}

\titledbox{Prediction Base Case}{
If we know the initial belief distribution $\probd{X_0}$ and want to compute the distribution of the next state $\probd{X_1}$ prior to any observations, we perform a 1-step prediction.

We can compute this by \linkterm{marginalizing}{def:marginalization} over all possible values of the initial state $X_0$:
\begin{align*}
\probd{X_1} &= \sum_{x_0 \in \text{dom}(X_0)} \probd{X_1, X_0 = x_0}
\end{align*}

Using the chain rule of probability, we split the joint probability into the \linkterm{transition model}{transition_model_hmm} and the initial \linkterm{prior}{def:bayesian_terms}:
\begin{align*}
\probd{X_1} &= \sum_{x_0} \probd{X_1 \mid X_0 = x_0} \cdot \probd{X_0 = x_0}
\end{align*}

This equation tells us that the probability of being in a state at $t=1$ is the sum of the probabilities of transitioning to it from every possible previous state, weighted by how likely those previous states were.
}

\minititle{Filtering Derivation}

To compute the belief state at time $t$ given all evidence up to time $t$, we want to find $\probd{X_t \mid E_{1:t}}$. We split the evidence into the current observation $E_t$ and all past observations $E_{1:t-1}$.

\textbf{Step 1: Update (Bayes' Rule / Product Rule)} \\
We want to find $P(X_t \mid E_t, E_{1:t-1})$. By definition of conditional probability:
\begin{align*}
\probd{X_t \mid E_{1:t}} &= \probd{X_t \mid E_t, E_{1:t-1}} \\
&= \frac{\probd{X_t, E_t, E_{1:t-1}}}{\probd{E_t, E_{1:t-1}}}
\end{align*}

We apply the product rule to the numerator ($\probd{A, B, C} = \probd{A \mid B, C} \cdot \probd{B, C}$):
\begin{align*}
\probd{X_t, E_t, E_{1:t-1}} &= \probd{E_t \mid X_t, E_{1:t-1}} \cdot \probd{X_t, E_{1:t-1}}
\end{align*}

And apply the product rule once more to $\probd{X_t, E_{1:t-1}}$:
\begin{align*}
\probd{X_t, E_t, E_{1:t-1}} &= \probd{E_t \mid X_t, E_{1:t-1}} \cdot \probd{X_t \mid E_{1:t-1}} \cdot \probd{E_{1:t-1}}
\end{align*}

We also apply the product rule to the denominator:
\begin{align*}
\probd{E_t, E_{1:t-1}} &= \probd{E_t \mid E_{1:t-1}} \cdot \probd{E_{1:t-1}}
\end{align*}

Now substitute both back into the fraction:
\begin{align*}
\probd{X_t \mid E_{1:t}} &= \frac{\probd{E_t \mid X_t, E_{1:t-1}} \cdot \probd{X_t \mid E_{1:t-1}} \cdot \probd{E_{1:t-1}}}{\probd{E_t \mid E_{1:t-1}} \cdot \probd{E_{1:t-1}}}
\end{align*}

The term $\probd{E_{1:t-1}}$ cancels out:
\begin{align*}
\probd{X_t \mid E_{1:t}} &= \frac{\probd{E_t \mid X_t, E_{1:t-1}} \cdot \probd{X_t \mid E_{1:t-1}}}{\probd{E_t \mid E_{1:t-1}}}
\end{align*}

Since $\frac{1}{\probd{E_t \mid E_{1:t-1}}}$ is just a normalizing constant, we can simplify our equation by replacing it with $\alpha$:
\begin{align*}
\probd{X_t \mid E_{1:t}} &= \alpha \Big( \probd{E_t \mid X_t, E_{1:t-1}} \cdot \probd{X_t \mid E_{1:t-1}} \Big)
\end{align*}

By the \linkterm{Sensor Markov property}{sensor_markov_property}, $E_t$ only depends on $X_t$, simplifying the first term (the \linkterm{sensor model}{sensor_model_hmm}):
\begin{align*}
\probd{X_t \mid E_{1:t}} &= \alpha \Big( \probd{E_t \mid X_t} \cdot \probd{X_t \mid E_{1:t-1}} \Big)
\end{align*}

\textbf{Step 2: Prediction (Law of Total Probability)} \\
Now we need to compute the 1-step prediction $\probd{X_t \mid E_{1:t-1}}$. We do this by \linkterm{marginalizing}{def:marginalization} over the previous state $X_{t-1}$:
\begin{align*}
\probd{X_t \mid E_{1:t-1}} &= \sum_{x} \probd{X_t, X_{t-1} = x \mid E_{1:t-1}}
\end{align*}

To show how the terms split, we expand the conditional joint probability using the definition of conditional probability:
\begin{align*}
\probd{X_t, X_{t-1} \mid E_{1:t-1}} &= \frac{\probd{X_t, X_{t-1}, E_{1:t-1}}}{\probd{E_{1:t-1}}}
\end{align*}

Applying the product rule to the numerator gives $\probd{X_t, X_{t-1}, E_{1:t-1}} = \probd{X_t \mid X_{t-1}, E_{1:t-1}} \cdot \probd{X_{t-1}, E_{1:t-1}}$:
\begin{align*}
\probd{X_t, X_{t-1} \mid E_{1:t-1}} &= \frac{\probd{X_t \mid X_{t-1}, E_{1:t-1}} \cdot \probd{X_{t-1}, E_{1:t-1}}}{\probd{E_{1:t-1}}}
\end{align*}

Notice that $\frac{\probd{X_{t-1}, E_{1:t-1}}}{\probd{E_{1:t-1}}}$ is simply the definition of conditional probability for $\probd{X_{t-1} \mid E_{1:t-1}}$. Therefore:
\begin{align*}
\probd{X_t, X_{t-1} \mid E_{1:t-1}} &= \probd{X_t \mid X_{t-1}, E_{1:t-1}} \cdot \probd{X_{t-1} \mid E_{1:t-1}}
\end{align*}

Substituting this conditionalized product rule back into our original marginalization sum gives:
\begin{align*}
\probd{X_t \mid E_{1:t-1}} &= \sum_{x} \probd{X_t \mid X_{t-1} = x, E_{1:t-1}} \cdot \probd{X_{t-1} = x \mid E_{1:t-1}}
\end{align*}

By the First-Order \linkterm{Markov property}{markov_chain}, the transition to $X_t$ only depends on $X_{t-1}$, simplifying the first term (the \linkterm{transition model}{transition_model_hmm}):
\begin{align*}
\probd{X_t \mid E_{1:t-1}} &= \sum_{x} \probd{X_t \mid X_{t-1} = x} \cdot \probd{X_{t-1} = x \mid E_{1:t-1}}
\end{align*}

\textbf{Step 3: Combine} \\
Substituting the prediction back into the update equation gives the recursive filtering formula. 
\begin{tcolorbox}[ams equation*,colback=yellow!10!white,colframe=red!50!black]
\probd{X_t \mid E_{1:t}^{=e}} = \alpha \left(
\probd{E_t = e_t \mid X_t} \cdot
\left(
\sum_{x} \probd{X_t \mid X_{t-1} = x} \cdot P(X_{t-1} = x \mid E_{1:t-1}^{=e})  
\right)
\right)
\end{tcolorbox}

\commandnote{Notice that $P(X_{t-1} = x \mid E_{1:t-1}^{=e})$ is simply the filtering result from the \textit{previous} time step. This means we can compute the belief state recursively!}


\anondefinition{
We call the inner part of the above expression the \defemph{forward} algorithm:
\[
\probd{X_t \mid E_{1:t}^{=e}} = \alpha \left(
\text{FORWARD} (e_t, \probd{X_{t-1} \mid E_{1:t-1}^{=e}})
\right) = f_{1:t}
\]
}{}

\centeredtitledbox{Filtering in Matrix Notation}{
\[
\probd{X_t \mid E_{1:t}^{=e}} = 
\alpha \left(
O_t \cdot T^T \cdot \probd{X_{t-1} \mid E_{1:t-1}^{=e}}
\right)
\]
}